\section{Evaluation targets and limits of interpretation}
\label{sec:final-evaluation-principles}

This thesis separates three questions that are often conflated in audio
authenticity research. First, does a measurement respond to its intended audio
phenomenon under controlled conditions? Second, does it discriminate the labels
in a specified collection? Third, does its utility survive a change in recording
source or generator? A positive answer to one question does not establish the
others. In particular, a classifier may exploit production or acquisition
conditions without identifying a consequence of generative synthesis.

\subsection{Labels and the unit of analysis}

The Human label denotes supported recording provenance, not an absence of MIDI,
sampling, quantization, pitch correction, compression, or editing. Conversely,
an AI-labelled recording may depend on human prompts, musical conditioning, or
post-production. Dynamics and timing statistics therefore cannot by themselves
establish whether a performance used MIDI or whether its authorship was human.

Recording identities, byte-identical audio objects, excerpts, derived stems,
and split groups are distinct units. Multiple views of one recording do not
provide independent evidence. Global groups preserve known dependencies across
sources and, where applicable, across labels. Cohort totals from different
duration experiments overlap and are not added to form a corpus size.
Unverified labels and pilot conditions are reported separately from eligible
classification populations.

\subsection{Duration and measurement availability}

Short recordings are not padded to create apparent evidence of long-term
structure. An unavailable beat sequence or an insufficient section observation
is not evidence of perfectly constant rhythm or zero structural variety.
Natural unobservability, processing failure, and an observed numerical zero are
retained as different states. A completed feature row is also not equivalent to
complete observation of every scalar in that row.

The exact-duration experiments control the physical analysis window within each
cohort. They do not isolate a causal effect of duration when their eligible
populations differ. Similarly, native-stereo eligibility narrows the population
to which stereo findings apply; duplicating mono channels does not manufacture
spatial information.

\subsection{Preprocessing, fitting, and decision scores}

Imputation and scaling are fitted on training observations only. The evaluation
retains the specified class, source, group, and recording weights, rather than
allowing large catalogues or duplicate views to dominate the objective. Training
quantity caps refer to global groups per source, not a universal number of
songs. Small sources are not duplicated to satisfy a nominal cap.

The frozen tabular classifier uses a ridge penalty of 10, an unpenalized
intercept, and a fixed decision threshold of 0.5. Its raw scores are unbounded
and are not calibrated posterior probabilities. Consequently, a score of 0.8
must not be described as an 80\% probability of AI generation. Missingness-only
and related diagnostics test selected alternative explanations; they do not
establish that every feature combination is free of processing artifacts.

\subsection{Separate evaluation directions}

Human-source-held-out evaluation examines transfer to an excluded Human source;
generator-held-out evaluation examines transfer to an excluded AI source or
declared generator family. Ordinary grouped evaluation is reported separately.
The latter can be informative about interpolation among observed sources while
remaining optimistic about an unseen generator. Results from these directions
are not interchangeable estimates of one deployment accuracy.

For studies using the composite criterion, let $A_H$ and $A_G$ denote the
specified source-macro AUC in the two transfer directions. Then
\[
  J = \tfrac{1}{2}(A_H + A_G).
\]
This is a mean of ranking statistics, not classification accuracy. Balanced
accuracy at the frozen threshold is reported separately, with the same declared
aggregation order. Scores produced by different fitted models are not casually
pooled into a cross-model AUC. Source-specific sensitivity and specificity remain
necessary even when a macro summary is favorable.

\subsection{Feature and training-quantity ablations}

Nonempty feature subsets and the five declared training caps are retained,
including negative increments. A useful standalone measurement can be redundant
or harmful in a larger model. Matched add-on contrasts therefore hold the cohort,
folds, quantity policy, preprocessing, weights, and threshold fixed wherever the
protocol specifies such a comparison. A maximum observed over many development
configurations is a descriptive result, not an independently validated winner.
Changes in sample count, weighting, and effective regularization also prevent a
training-cap contrast from being interpreted as a pure causal effect of data
diversity.

\subsection{Uncertainty and evidence scope}

Conditional resampling of held-out recording groups describes uncertainty
within fixed fitted models and observed recordings. It does not include
retraining or model-selection variability. Separate resampling of observed source
labels characterizes sensitivity to that source composition; it does not turn
a purposively assembled collection into an independent sample of future music
domains. Undefined cells and the declared aggregation policy are retained rather
than replaced by zero or chance performance.

Software tests, byte hashes, schema audits, independent numerical replay,
controlled interventions, and external classification answer different questions.
A hash confirms identity within its stated scope; it does not prove scientific
validity. Numerical replay verifies selected arithmetic, not a causal account of
authorship. A controlled phase-coupling response, for example, establishes a
measurement capability under those conditions, not an AI/Human classification
result. Failed admission gates and original failed audits are retained alongside
any justified implementation corrections.

\subsection{Reproducibility references}

The authoritative implementations and result-specific protocols are indexed in
the reproducibility appendix. Relevant preserved implementations include
\path{extract_native30_sdrp_operational_v1.py},
\path{evaluate_native30_v3.py}, and
\path{bootstrap_sdrfh10_uncertainty.py}. Their existence is not asserted as
proof that an associated production run completed; empirical chapters require
the corresponding committed outputs and verification evidence.
